\documentclass[10pt]{article}
\usepackage[T1]{fontenc}
\usepackage[utf8]{inputenc}
\usepackage[french]{babel}
\usepackage{amsmath,amssymb,mathrsfs,mathpazo}
\usepackage{geometry}
\geometry{a4paper,margin=2cm}
\usepackage{graphicx,array,tabularx,tabularray,colortbl,multirow,multicol,booktabs,enumitem,pifont,ulem,fancybox,diagbox,cancel}
\usepackage{pstricks,pstricks-add,pst-plot,pst-tree,pst-node,pst-text,pst-func,pst-math,pst-bspline,pst-3dplot,pst-eucl,pst-solides3d}
\let\pscalculate\undefined
\usepackage{pst-bezier}
\usepackage{tikz,pgfplots,tkz-tab}
\usetikzlibrary{arrows,arrows.meta,patterns,calc,positioning,trees,decorations.markings,shapes.geometric,shapes.misc,shapes.arrows,intersections,decorations.pathmorphing,decorations.pathreplacing,backgrounds,shadows,babel}
\pgfplotsset{compat=1.18}
\usepgfplotslibrary{fillbetween}
\usepackage[np]{numprint}
\usepackage{listings,makecell,pas-tableur,twoopt}
\let\barre\undefined
\usepackage[pstricks]{tablvar}
\definecolor{darkgreen}{RGB}{0,100,0}
\providecommand{\red}{\color{red}}
\providecommand{\blue}{\color{blue}}
\providecommand{\green}{\color{green}}
\providecommand{\black}{\color{black}}
\providecommand{\white}{\color{white}}
\providecommand{\gray}{\color{gray}}
\providecommand{\cyan}{\color{cyan}}
\providecommand{\magenta}{\color{magenta}}
\providecommand{\yellow}{\color{yellow}}
\providecommand{\violet}{\color{violet}}
\providecommand{\orange}{\color{orange}}
\providecommand{\euro}{\texteuro}
\providecommand{\EUR}{\texteuro}
\providecommand{\ieme}{\textsuperscript{e}}
\providecommand{\degres}{\ensuremath{^\circ}}
\providecommand{\degree}{\ensuremath{^\circ}}
\providecommand{\vv}[1]{\overrightarrow{#1}}
\providecommand{\Coord}[1]{#1}
\providecommand{\decosix}{}
\providecommand{\starredbullet}{\ensuremath{\bullet}}
\pagestyle{empty}
\setlength{\parindent}{0pt}
\usepackage{adjustbox}\DeclareUnicodeCharacter{2236}{:}\DeclareUnicodeCharacter{2190}{\ensuremath{\leftarrow}}
\providecommand{\R}{\mathbb{R}}
\providecommand{\N}{\mathbb{N}}
\providecommand{\Z}{\mathbb{Z}}
\providecommand{\Q}{\mathbb{Q}}
\providecommand{\D}{\mathbb{D}}
\providecommand{\C}{\mathbb{C}}
\providecommand{\e}{{\mathrm{e}}}
\providecommand{\vect}[1]{\overrightarrow{#1}}
\providecommand{\vectt}[1]{\overrightarrow{\mathrm{#1}}}
\providecommand{\vvt}[1]{\overrightarrow{\mathrm{#1}}}
\providecommand{\np}[2][]{{#2~#1}}
\providecommand{\vv}[1]{{\overrightarrow{#1}}}
\renewcommand{\d}{{\mathrm{d}}}
\providecommand{\degres}{{^\circ}}
\providecommand{\degre}{{^\circ}}
\providecommand{\blue}{\color{blue}}
\providecommand{\red}{\color{red}}
\providecommand{\enskip}{\quad}
\providecommand{\llbracket}{\lbrack\!\lbrack}
\providecommand{\rrbracket}{\rbrack\!\rbrack}
\providecommand{\euro}{{\text{€}}}
\providecommand{\no}{{\text{n°}}}
\providecommand{\interval}[2]{{[#1;#2]}}
\providecommand{\intervFO}[2]{{[#1;#2[}}
\providecommand{\dotfill}{\ldots}

\providecommand{\ds}{\displaystyle}
\providecommand{\textding}[1]{#1}
\providecommand{\Oij}{\(\mathrm{O};\vec{\imath},\vec{\jmath}\)}
\providecommand{\Oijk}{\(\mathrm{O};\vec{\imath},\vec{\jmath},\vec{k}\)}
\providecommand{\Ouv}{\(\mathrm{O};\vec{u},\vec{v}\)}
\providecommand{\intervFF}[2]{\ensuremath{[#1;#2]}}
\renewcommand{\euro}{\ifmmode\text{\texteuro}\else\texteuro\fi}
\DeclareUnicodeCharacter{00B0}{\ifmmode^{\circ}\else\textdegree\fi}
\DeclareUnicodeCharacter{20AC}{\ifmmode\text{\texteuro}\else\texteuro\fi}
% Énoncés : Éducation nationale. Transcription et corrigés : APMEP, Denis Vergès et contributeurs. Découpage et classement : CoopMaths, Sébastien Lozano, Éric Elter et contributeurs. Adaptation de la présentation : Maths Explorer.
% CC BY-SA : https://coopmaths.fr/www/about/
% https://coopmaths.fr/alea/static/eam/2026/tex/eam_technologique_2026_06_metropole_3_cor.tex
\begin{document}
\noindent\textit{APMEP / CoopMaths. CC BY-SA. Présentation adaptée par Maths Explorer.}\par\medskip
\par
\textbf{Pour chaque affirmation, répondre par VRAI ou FAUX, en justifiant la réponse.}\par
Toute absence de justification ou justification incorrecte ne sera pas prise en compte dans la notation.

Une enquête est réalisée sur l'équipement informatique des 400 élèves d'un collège.
\begin{center}\renewcommand{\arraystretch}{1.35}
\begin{tabular}{|p{3.2cm}|c|c|c|}\hline
& Ont une adresse mail & N'ont pas d'adresse mail & Total\\\hline
Ont un équipement individuel & 50 & 210 & 260\\\hline
N'ont pas d'équipement individuel & 40 & 100 & 140\\\hline
Total & 90 & 310 & 400\\\hline
\end{tabular}
\end{center}

\begin{enumerate}
\item \underline{Affirmation}

50\% des élèves du collège possèdent une adresse mail et un équipement individuel.

\item \underline{Affirmation}

Au moins 50\% des élèves du collège ne possèdent pas d'adresse mail.

\item \underline{Affirmation}

On choisit un élève au hasard dans le collège. La probabilité qu'il ne possède ni adresse mail ni équipement individuel est égale à 25\%.

\item \underline{Affirmation}

Au moins $\dfrac{1}{5}$ des élèves ayant un équipement individuel possèdent également une adresse mail.
\end{enumerate}

\begin{color}{blue}Solution :\par

\begin{enumerate}
\item
50 élèves du collège possèdent une adresse mail et un équipement individuel sur 400 élèves au total. Cela correspond à une proportion de $\dfrac{50}{400} = \dfrac{12,5}{100} = 12,5~\% \neq 50~\% $.

\textbf{L'affirmation est fausse.}

\item 
310 élèves du collège ne possèdent pas d'adresse mail. $310>200$, la moitié de 400.

\textbf{L'affirmation est vraie.}
\item 

100 élèves du collège ne possèdent ni d'adresse mail ni d'équipement individuel sur 400 élèves au total. Cela correspond à une proportion de $\dfrac{100}{400} = \dfrac{25}{100} = 25~\% $. Si on choisit un élève au hasard parmi les élèves du collège, avec l'hypothèse d'équiprobabilité, la probabilité qu'il ne possède ni adresse mail ni équipement individuel est égale à 25\%.

\textbf{L'affirmation est vraie.}
\item 

260 élèves possèdent un équipement individuel. Parmi ceux-ci, 50 ont une adresse mail. La proportion d'élèves ayant une adresse mail parmi les élèves ayant un équipement individuel est $\dfrac{50}{260}=\dfrac{5}{26}$ et $\dfrac{5}{26}<\dfrac{5}{25}$. Comme $\dfrac{5}{25}=\dfrac{1}{5}$, $\dfrac{50}{260}<\dfrac{1}{5}$.

\textbf{L'affirmation est fausse.}

\end{enumerate}

\end{color}

\end{document}
